\subsection{Filtered rings and modules}

DEFINITION 1. An increasing (resp. decreasing) sequence \((\mathbf{G}_n)_{n\in \mathbf{Z}}\) of subgroups of a group \(\mathbf{G}\) is called an increasing (resp. decreasing) filtration on \(\mathbf{G}\) .

A group with a filtration is called a filtered group.

If \((\mathrm{G}_{n})_{n\in\mathbf{Z}}\) is an increasing (resp. decreasing) filtration on a group G and we write \(G_{n}^{\prime}=G_{-n}\) , clearly \((\mathrm{G}_{n}^{\prime})_{n\in\mathbf{Z}}\) is a decreasing (resp. increasing) filtration on G. We may therefore restrict our study to decreasing filtrations and henceforth when we speak of a filtration, we shall mean a decreasing filtration, unless otherwise stated.

Given a decreasing filtration \((\mathbf{G}_n)_{n\in \mathbf{Z}}\) on a group \(\mathbf{G}\) , clearly \(\bigcap_{n\in \mathbf{Z}}\mathbf{G}_n\) and \(\bigcup_{n\in \mathbf{Z}}\mathbf{G}_n\) are subgroups of \(\mathbf{G}\) ; the filtration is called separated if \(\bigcap_{n\in \mathbf{Z}}\mathbf{G}_n\) is reduced to the identity element and exhaustive if \(\bigcup_{n\in \mathbf{Z}}\mathbf{G}_n = \mathbf{G}\) .

DEFINITION 2. Given a ring A, a filtration \((\mathbf{A}_n)_{n\in \mathbf{Z}}\) over the additive group A is called compatible with the ring structure on A if

\[
\mathrm{A} _ {m} \mathrm{A} _ {n} \subset \mathrm{A} _ {m + n} \quad for \quad m \in \mathbf {Z}, \quad n \in \mathbf {Z}\tag{1}
\]

\[
\mathbf {1} \in \mathbf {A} _ {0}.\tag{2}
\]

The ring A with this filtration is then called a filtered ring.

Conditions (1) and (2) show that \(\mathbf{A}_0\) is a subring of \(\mathbf{A}\) and the \(\mathbf{A}_n\) (left and right) \(\mathbf{A}_0\) -modules. The set \(\mathbf{B} = \bigcup_{n \in \mathbf{Z}} \mathbf{A}_n\) is a subring of \(\mathbf{A}\) and the set \(\mathfrak{n} = \bigcap_{n \in \mathbf{Z}} \mathbf{A}_n\) a two-sided ideal of \(\mathbf{B}\) ; for if \(x \in \mathfrak{n}\) and \(a \in \mathbf{A}_p\) , for all \(k \in \mathbf{Z}\) , \(x \in \mathbf{A}_{k-p}\) , whence \(ax \in \mathbf{A}_k\) and \(xa \in \mathbf{A}_k\) by (1); therefore \(ax \in \mathfrak{n}\) and \(xa \in \mathfrak{n}\) .

An important particular case is that in which \(\mathbf{A}_0 = \mathbf{A}\) ; then \(\mathbf{A}_n = \mathbf{A}\) for \(n \leqslant 0\) and all the \(\mathbf{A}_n\) are two-sided ideals of \(\mathbf{A}\) .

DEFINITION 3. Let A be a filtered ring, \((\mathbf{A}_n)_{n\in \mathbb{Z}}\) its filtration and E an A-module. A filtration \((\mathrm{E}_n)_{n\in \mathbb{Z}}\) on E is called compatible with its module structure over the filtered ring A if

\[
\mathrm{A} _ {m} \mathrm{E} _ {n} \subset \mathrm{E} _ {m + n} \quad for \quad m \in \mathbf {Z}, \quad n \in \mathbf {Z}.\tag{3}
\]

The A-module E with this filtration is called a filtered module.

The \(\mathbf{E}_n\) are all \(\mathbf{A}_0\) -modules; if \(\mathbf{B} = \bigcup_{n \in \mathbf{Z}} \mathbf{A}_n\) , clearly \(\bigcup_{n \in \mathbf{Z}} \mathbf{E}_n\) is a B-module and so is \(\bigcap_{n \in \mathbf{Z}} \mathbf{E}_n\) by the same argument as above for \(\bigcap_{n \in \mathbf{Z}} \mathbf{A}_n\) . If \(\mathbf{A}_0 = \mathbf{A}\) , all the \(\mathbf{E}_n\) are submodules of \(\mathbf{E}\) .

\textbf{Examples}\par

\begin{enumerate}
\def\labelenumi{(\arabic{enumi})}
\tightlist
\item
  Let A be a graded ring of type Z; for all \(i \in Z\) , let \(\mathbf{A}_{(i)}\) be the subgroup of homogeneous elements of degree i in A. Let us write \(A_{n} = \sum_{i \geq n} A_{(i)}\) ; then it is immediate that \((\mathbf{A}_{n})\) is an exhaustive and separated decreasing filtration which is compatible with the ring structure on A; this filtration is said to be associated with the graduation \((\mathbf{A}_{(i)})_{i \in \mathbf{Z}}\) and the filtered ring A is said to be associated with the given graded ring A.
\end{enumerate}

Now let E be a graded module of type Z over the graded ring A and for all \(i \in Z\) let \(\mathrm{E}_{(i)}\) be the subgroup of homogeneous elements of degree i of E. Let us write \(\mathrm{E}_{n} = \sum_{i \geqslant n} \mathrm{E}_{(i)}\) ; then \((\mathrm{E}_{n})\) is an exhaustive and separated decreasing filtration which is compatible with the module structure on E over the filtered ring A; this filtration is said to be associated with the graduation \((\mathrm{E}_{(i)})_{i \in \mathbf{Z}}\) and the filtered module E is said to be associated with the given graded module E. (2) Let A be a filtered ring, \((\mathrm{A}_{n})_{n \in \mathbf{Z}}\) its filtration and E an A-module. Let us write \(E_{n} = A_{n}E\) ; it follows from (1) that

\[
\mathbf {A} _ {m} \mathbf {E} _ {n} = \mathbf {A} _ {m} \mathbf {A} _ {n} \mathbf {E} \subset \mathbf {A} _ {m + n} \mathbf {E} = \mathbf {E} _ {m + n},
\]

and from (2) that \(E_{0} = E\) ; then \((\mathrm{E}_{n})\) is an exhaustive filtration which is compatible with the A-module structure on E. This filtration is said to be derived from the given filtration \((\mathrm{A}_{n})\) on A; note that it is not necessarily separated, even if \((\mathbf{A}_n)\) is separated and \(\mathbf{E}\) and the \(\mathbf{A}_n\) are finitely generated A-modules (cf.~§ 3, Exercise 2; see however § 3, no. 3, Proposition 5 and no. 2, Corollary to Proposition 4).

\begin{enumerate}
\def\labelenumi{(\arabic{enumi})}
\setcounter{enumi}{2}
\tightlist
\item
  Let A be a ring and m a two-sided ideal of A. Let us write \(A_{n} = m^{n}\) for \(n \geqslant 0\) , \(A_{n} = A\) for n \textless{} 0; it is immediate that \((\mathbf{A}_{n})\) is an exhaustive filtration on A, called the m-adic filtration. Let E be an A-module; the filtration \((\mathbf{E}_{n})\) derived from the m-adic filtration on A is called the m-adic filtration on E; in other words, \(E_{n} = m^{n}E\) for \(n \geqslant 0\) and \(E_{n} = E\) for n \textless{} 0.
\end{enumerate}

If A is commutative and B is an A-algebra, n = mB is a two-sided ideal of B and, for every B-module F, \(n^{k}F = m^{k}F\) and hence the n-adic filtration on F coincides with the m-adic filtration (if F is considered as an A-module).

\begin{enumerate}
\def\labelenumi{(\arabic{enumi})}
\setcounter{enumi}{3}
\item
  If A is a filtered ring and \((\mathbf{A}_{n})\) its filtration, the left A-module \(A_{s}\) is a filtered A-module with the filtration \((\mathbf{A}_{n})\) . On the other hand, clearly \((\mathbf{A}_{n})\) is a filtration which is compatible with the ring structure on the opposite ring \(A^{0}\) and \(A_{d}\) is a filtered (left) \(A^{0}\) -module with the filtration \((\mathbf{A}_{n})\) .
\item
  On a ring A the sets \(A_{n}\) such that \(A_{n}=0\) for n\textgreater0, \(A_{n}=A\) for \(n\leqslant0\) form what is called a trivial filtration associated (Example 1) with the trivial graduation on A; on an A-module E, every filtration ( \(E_{n}\) ) consisting of sub-A-modules is then compatible with the module structure on E over the filtered ring A. Then it is possible to say that every filtered commutative group G is a filtered Z-module, if Z is given the trivial filtration.
\end{enumerate}

Let \(G\) be a filtered group and \((G_n)_{n \in \mathbf{Z}}\) its filtration; clearly, for every subgroup \(H\) of \(G\) , \((H \cap G_n)_{n \in \mathbf{Z}}\) is a filtration said to be induced by that on \(G\) ; it is exhaustive (resp. separated) if that on \(G\) is. Similarly, if \(H\) is a normal subgroup of \(G\) , the family \(((H.G_n)/H)_{n \in \mathbf{Z}}\) is a filtration on the group \(G/H\) , called the quotient under \(H\) of the filtration on \(G\) ; it is exhaustive if \((G_n)\) is. If \(G'\) is another filtered group and \((G_n')_{n \in \mathbf{Z}}\) its filtration, \((G_n \times G_n')\) is a filtration on \(G \times G'\) called the product of the filtrations on \(G\) and \(G'\) , which is exhaustive (resp. separated) if \((G_n)\) and \((G_n')\) are.

Now let A be a filtered ring and \((\mathbf{A}_{n})\) its filtration; on every subring B of A, clearly the filtration induced by that on A is compatible with the ring structure on B. If b is a two-sided ideal of A, the quotient filtration on A/b of that on A is compatible with the structure of this ring, for

\[
\left(\mathrm{A} _ {n} + \mathfrak {b}\right) \left(\mathrm{A} _ {m} + \mathfrak {b}\right) \subset \mathrm{A} _ {n + m} + \mathfrak {b}.
\]

If \(A'\) is another filtered ring, the product filtration on \(A \times A'\) is compatible with the structure of this ring.

Finally let E be a filtered A-module and \((E_{n})\) its filtration; on every sub-module F of E, the filtration induced by that on E is compatible with the A-module structure on F and, on the quotient module E/F, the quotient filtration of that on E is compatible with the A-module structure, as

\[
\mathrm{A} _ {n} (\mathrm{F} + \mathrm{E} _ {m}) \subset \mathrm{F} + \mathrm{A} _ {n} \mathrm{E} _ {m} \subset \mathrm{F} + \mathrm{E} _ {m + n}.
\]

\pandocbounded{\includegraphics[keepaspectratio]{images/part002_676c5d74.jpg}}

Note that if the filtration on E is derived from that on A (Example 2), so is the quotient filtration on E/F, but not in general the filtration induced on F (Exercise 1; see however § 3, no. 2, Theorem 2).

If \(E'\) is another filtered A-module, the product filtration on \(E \times E'\) is compatible with its A-module structure. If the filtrations on E and \(E'\) are derived from that on A (Example 2), so is their product filtration.
% semantic-fragments: legacy-input-seam
\relax{}
